% !PRINTED 342
\weberpage{342}

Man kann die Bezeichnung der vier Systeme von geraden
Linien auch aus dem einen Schema:
\[
\begin{array}{ccc}
(0,\,0) & (1,\,2) & (2,\,1)\\
(1,\,1) & (2,\,0) & (0,\,2)\\
(2,\,2) & (0,\,1) & (1,\,0)
\end{array}
\]
ableiten, wenn man die Zeilen, die Colonnen und die je drei
den positiven und negativen Gliedern der Determinantenbildung
entsprechenden Combinationen aufstellt.

\S.\,92.

\gesperrt{Tripelgleichungen.}

Das System der neun Wendepunkte hat die Eigenschaft,
dass man aus zwei beliebigen von ihnen einen ganz bestimmten
dritten auf rationalem Wege ableiten kann. Dies drückt sich
als eine Eigenschaft der Gleichung $9^{\text{ten}}$ Grades aus, von der die
Bestimmung der Wendepunkte abhängt, die z.\,B. die Abscissen
dieser Punkte zu Wurzeln hat. Wir stellen folgende Definition
auf:

\gesperrt{Eine Gleichung ohne gleiche Wurzeln heisst eine
Tripelgleichung, wenn je zwei ihrer Wurzeln eine
dritte bestimmen, so dass jede Wurzel eines solchen
Tripels rational durch die beiden anderen ausgedrückt
werden kann, und zwar so, dass man in einer Gleichung}
\weq{1}{x_3 = \Theta\,(x_1, x_2),}
\gesperrt{in der $\Theta$ eine festgehaltene rationale Function bedeutet, für
$x_1, x_2, x_3$ die Wurzeln irgend eines Tripels
in beliebiger Reihenfolge setzen kann}\footnote{Die Tripelgleichungen $9^{\text{ten}}$ Grades sind zuerst behandelt von
\gesperrt{Hesse} in der Abhandlung: \glqq Algebraische Auflösung derjenigen Gleichung
$9^{\text{ten}}$ Grades etc.\grqq{} in Crelle's Journal, Bd.\,34 (1847).}.

Die Gleichung, von der die Wendepunkte abhängen, ist eine
Tripelgleichung $9^{\text{ten}}$ Grades. Es giebt aber auch Tripelgleichungen
von anderem, z.\,B. vom $7^{\text{ten}}$ Grade, zu denen die im vorigen
Abschnitte betrachteten gehören\footnote{Vgl. \gesperrt{N\"other}, Mathem. Annalen, Bd.\,15 (1879): \glqq Ueber die Gleichungen
$8^{\text{ten}}$ Grades und ihr Auftreten in der Theorie der Curven vierter
Ordnung.\grqq{}}.
